% =============================================================================
% PARTE II: TOXICOLOGÍA DE GASES, DINÁMICA DE HUMOS Y MARCO LEGAL ARGENTINO
% =============================================================================

% --- SLIDE 16: MONÓXIDO DE CARBONO ---
\begin{frame}{Toxicología de Gases Asfixiantes: Monóxido de Carbono ($CO$)}
  \begin{columns}[c]
    \begin{column}{0.50\textwidth}
      \small
      El $CO$ es el causante primario de letalidad en siniestros estructurales:
      \begin{itemize}
        \scriptsize
        \item Gas inodoro, incoloro e insípido; no advierte sensorialmente.
        \item Posee una afinidad química por el hierro de la hemoglobina férrica aproximadamente \textbf{250 veces superior} a la del oxígeno.
        \item Forma \textbf{Carboxihemoglobina ($COHb$)}, desplazando al $O_2$ y bloqueando su entrega a los tejidos.
        \item \textbf{Efecto físico fulminante:} Pérdida rápida de agudeza visual, confusión y debilidad muscular en miembros inferiores que impide la bipedestación.
      \end{itemize}
    \end{column}

    \begin{column}{0.48\textwidth}
      \centering
      \resizebox{0.9\columnwidth}{!}{%
        \begin{tikzpicture}
          \node[draw=utnblue, fill=utnblue!10, rounded corners=6pt, line width=1.2pt, inner sep=6pt, text width=4.5cm, align=center] (Hb) {
            \textbf{Molécula de Hemoglobina}\\[2pt]
            \scriptsize Capacidad de transporte bloqueada
          };
          \node[draw=firered, fill=firered!20, circle, line width=1.2pt, font=\scriptsize\bfseries, above=0.6cm of Hb] (CO) {CO};
          \draw[->, >=stealth, line width=2pt, firered] (CO) -- node[right, font=\tiny\bfseries, firered] {Afinidad $\times 250$} (Hb);
        \end{tikzpicture}
      }
    \end{column}
  \end{columns}
\end{frame}

% --- SLIDE 17: CIANURO DE HIDRÓGENO ---
\begin{frame}{Toxicología Celular: Cianuro de Hidrógeno ($HCN$)}
  \begin{columns}[c]
    \begin{column}{0.50\textwidth}
      \small
      El $HCN$ actúa como un veneno respiratorio a escala intracelular:
      \begin{itemize}
        \scriptsize
        \item Se desprende durante la descomposición térmica de materiales que contienen nitrógeno: espumas de poliuretano de aislamiento, tapizados, nylon y lanas.
        \item Difunde a través de la membrana plasmática e \textbf{inhibe irreversiblemente la enzima citocromo c oxidasa} en la mitocondria.
        \item La célula se vuelve incapaz de captar y utilizar el oxígeno transportado por la sangre, deteniendo la síntesis de ATP.
        \item \textbf{Sinergia letal con $CO$:} El $CO$ bloquea la entrega y el $HCN$ bloquea la absorción celular.
      \end{itemize}
    \end{column}

    \begin{column}{0.48\textwidth}
      \centering
      \resizebox{0.9\columnwidth}{!}{%
        \begin{tikzpicture}
          \node[draw=darkcharcoal, fill=cardbg, rounded corners=6pt, line width=1.2pt, inner sep=8pt, text width=4.5cm, align=center] {
            \textbf{\textcolor{firered}{ANPOXIA HISTOTÓXICA}}\\[4pt]
            \scriptsize Colapso de la cadena de fosforilación oxidativa mitocondrial.\\[3pt]
            \scriptsize Pérdida de conciencia en menos de 2 a 3 minutos.
          };
        \end{tikzpicture}
      }
    \end{column}
  \end{columns}
\end{frame}

% --- SLIDE 18: GASES ÁCIDOS E IRRITANTES ---
\begin{frame}{Gases Ácidos, Corrosivos e Irritantes de la Combustión}
  \begin{columns}[T]
    \begin{column}{0.48\textwidth}
      \begin{block}{\textbf{Cloruro de Hidrógeno ($HCl$)}}
        \small
        $\bullet$ Generado en la pirólisis de aislamientos de cables de \textbf{PVC}.\\
        $\bullet$ En contacto con el vapor de agua y las mucosas forma ácido clorhídrico concentrado.\\
        $\bullet$ Provoca quemaduras químicas inmediatas en tráquea y alvéolos.\\
        $\bullet$ \textbf{Impacto electrónico:} Corroe pistas de cobre y uniones de silicio de circuitos impresos vecinos.
      \end{block}
    \end{column}

    \begin{column}{0.48\textwidth}
      \begin{block}{\textbf{Óxidos de Nitrógeno y Azufre}}
        \small
        $\bullet$ \textbf{Dióxido de nitrógeno ($NO_2$):} Desencadena edema pulmonar diferido a las 24 horas del siniestro.\\
        $\bullet$ \textbf{Dióxido de azufre ($SO_2$) y acroleína:} Irritantes sensoriales agudos que provocan \textbf{blefaroespasmo involuntario} (cierre de ojos forzado) y anulación visual del escape.
      \end{block}
    \end{column}
  \end{columns}
\end{frame}

% --- SLIDE 19: CONCENTRACIÓN DE OXÍGENO ---
\begin{frame}{Umbrales Fisiológicos de Caída del Oxígeno Ambiental}
  \small
  La combustión en recintos cerrados consume vorazmente el comburente disponible en el aire:

  \vspace{0.2cm}
  \centering
  \resizebox{0.95\textwidth}{!}{%
    \begin{tikzpicture}[yscale=0.75, line cap=round, line join=round]
      % Escala vertical
      \draw[line width=1.5pt, ->, >=stealth, darkcharcoal] (0,0) -- (0,5.8) node[above, font=\scriptsize\bfseries] {Concentración de $O_2$};

      % Niveles
      \draw[fill=successgreen!20, draw=successgreen!80!black, line width=1pt] (0.5,4.5) rectangle (9.5,5.3) 
            node[midway, font=\scriptsize] {\textbf{20.9\% (Normal):} Respiración y funciones psicomotrices plenamente estables.};

      \draw[fill=amberyellow!25, draw=amberyellow!90!black, line width=1pt] (0.5,3.3) rectangle (9.5,4.1) 
            node[midway, font=\scriptsize] {\textbf{16.0\% (Umbral de Alerta):} Taquicardia refleja, desorientación espacial y merma de juicio.};

      \draw[fill=fireorange!25, draw=fireorange!90!black, line width=1pt] (0.5,2.1) rectangle (9.5,2.9) 
            node[midway, font=\scriptsize] {\textbf{14.0\% (Falla Motriz):} Pérdida de fuerza muscular en piernas; el operador cae al piso.};

      \draw[fill=firered!25, draw=firered!90!black, line width=1pt] (0.5,0.9) rectangle (9.5,1.7) 
            node[midway, font=\scriptsize] {\textbf{10.0\% (Crítico):} Náuseas violentas, vómitos, pérdida de reflejos y coma profundo.};

      \draw[fill=darkcharcoal!30, draw=darkcharcoal, line width=1.2pt] (0.5,-0.3) rectangle (9.5,0.5) 
            node[midway, font=\scriptsize\bfseries, firered] {\textbf{$<$ 6.0\% (Letal Fulminante):} Cese total de respiración espontánea en $< 40\,\mathrm{s}$.};
    \end{tikzpicture}
  }
\end{frame}

% --- SLIDE 20: DINÁMICA DEL HUMO Y PLANO NEUTRO ---
\begin{frame}{Dinámica del Humo y el Concepto de Plano Neutro}
  \begin{columns}[c]
    \begin{column}{0.48\textwidth}
      \small
      Evolución espacial del penacho térmico en un recinto:
      \begin{itemize}
        \scriptsize
        \item Los humos sobrecalentados ascienden por flotabilidad y chocan con el techo (\textit{ceiling jet}).
        \item Se acumulan en una \textbf{capa superior de gases calientes y hollín}.
        \item Se genera un \textbf{Plano Neutro}: límite dinámico de separación entre la masa de humo caliente superior y el aire limpio respirable inferior.
        \item \textbf{Criterio de Evacuación:} Desplazamiento gateando o agachado por debajo del plano neutro para conservar visión y oxígeno.
      \end{itemize}
    \end{column}

    \begin{column}{0.50\textwidth}
      \centering
      \resizebox{0.95\columnwidth}{!}{%
        \begin{tikzpicture}[scale=0.9]
          % Habitación
          \draw[line width=1.5pt, darkcharcoal] (0,0) rectangle (5,3.2);

          % Capa superior de humo
          \fill[darkcharcoal!60] (0,1.8) rectangle (5,3.2);
          \node[font=\scriptsize\bfseries, white] at (2.5,2.5) {Humo y Gases Calientes ($>200\,\Celsius$)};

          % Plano Neutro
          \draw[dashed, line width=1.5pt, fireorange] (0,1.8) -- (5,1.8) node[right, font=\tiny\bfseries, fireorange] {Plano Neutro};

          % Aire limpio inferior
          \node[font=\scriptsize\bfseries, successgreen!80!black] at (2.5,0.8) {Aire Frío Respirable ($O_2 \approx 20\%$)};

          % Foco
          \node[circle, fill=firered, inner sep=4pt] at (0.8,0.3) {};
          \draw[->, >=stealth, line width=1.2pt, firered] (0.8,0.4) .. controls (1.0,1.2) .. (1.5,2.2);
        \end{tikzpicture}
      }
    \end{column}
  \end{columns}
\end{frame}

% --- SLIDE 21: DIAGNÓSTICO VISUAL DE HUMOS ---
\begin{frame}{Diagnóstico Pericial del Humo según su Coloración}
  \small
  El color del humo revela de manera inmediata la química del combustible y el nivel de ventilación:

  \vspace{0.2cm}
  \centering
  \resizebox{0.98\textwidth}{!}{%
    \begin{tikzpicture}[yscale=0.7]
      \draw[fill=white, draw=grayborder, line width=1pt] (0,4) rectangle (10,4.8)
            node[midway, font=\scriptsize] {\textbf{Blanco:} Combustión completa de celulosa seca / Presencia abundante de vapor de agua.};

      \draw[fill=amberyellow!30, draw=amberyellow!90!black, line width=1pt] (0,3) rectangle (10,3.8)
            node[midway, font=\scriptsize] {\textbf{Amarillo:} Sustancias químicas con azufre o ácidos nítrico/clorhídrico; gases letales.};

      \draw[fill=gray!30, draw=gray, line width=1pt] (0,2) rectangle (10,2.8)
            node[midway, font=\scriptsize] {\textbf{Gris:} Celulosa, paja y tejidos en combustión incompleta con deficiencia de aire.};

      \draw[fill=darkcharcoal!35, draw=darkcharcoal!80, line width=1pt] (0,1) rectangle (10,1.8)
            node[midway, font=\scriptsize] {\textbf{Negro Claro:} Cauchos sintéticos, carbón mineral y polímeros de cadena simple.};

      \draw[fill=darkcharcoal!85, draw=black, line width=1pt] (0,0) rectangle (10,0.8)
            node[midway, font=\scriptsize, white] {\textbf{Negro Oscuro:} Hidrocarburos pesados, naftas, plásticos espumados y fibras acrílicas.};
    \end{tikzpicture}
  }
\end{frame}

% --- SLIDE 22: MARCO LEGAL ARGENTINO ---
\begin{frame}{Estructura del Marco Regulatorio Argentino}
  \begin{columns}[c]
    \begin{column}{0.48\textwidth}
      \small
      El ordenamiento jurídico en higiene y seguridad industrial se asienta en tres niveles normativos:
      \begin{enumerate}
        \scriptsize
        \item \textbf{Ley Nacional 19.587 (1972):} Marco general obligatorio de condiciones de higiene y seguridad en todo el territorio argentino.
        \item \textbf{Ley 24.557 (1995):} Régimen de Prevención y Reparación de Riesgos del Trabajo (SRT y ART).
        \item \textbf{Decreto Reglamentario 351/79:}
        \begin{itemize}
          \tiny
          \item \textbf{Capítulo 18:} Protección activa, pasiva y extintores.
          \item \textbf{Anexo VII:} Cálculos de carga de fuego y medios de escape.
          \item \textbf{Capítulo 21:} Capacitación obligatoria registrada.
        \end{itemize}
      \end{enumerate}
    \end{column}

    \begin{column}{0.48\textwidth}
      \centering
      \resizebox{0.85\columnwidth}{!}{%
        \begin{tikzpicture}[scale=0.9]
          \draw[fill=utnblue!90, draw=black, line width=1pt] (0,3.2) rectangle (4,4.2) node[midway, white, font=\scriptsize\bfseries] {Ley 19.587 / Ley 24.557};
          \draw[fill=utncyan!80, draw=black, line width=1pt] (0,1.6) rectangle (4,2.6) node[midway, white, font=\scriptsize\bfseries] {Decreto 351/79 (Anexo VII)};
          \draw[fill=lightbg!80!black, draw=black, line width=1pt] (0,0) rectangle (4,1.0) node[midway, darkcharcoal, font=\scriptsize\bfseries] {Normas IRAM / AEA / NFPA};

          \draw[->, >=stealth, line width=1.5pt, firered] (2,3.2) -- (2,2.6);
          \draw[->, >=stealth, line width=1.5pt, firered] (2,1.6) -- (2,1.0);
        \end{tikzpicture}
      }
    \end{column}
  \end{columns}
\end{frame}

% --- SLIDE 23: CARGA DE FUEGO FÓRMULA ---
\begin{frame}{Determinación Analítica de la Carga de Fuego ($Q_f$)}
  \begin{columns}[c]
    \begin{column}{0.50\textwidth}
      \begin{block}{\textbf{Fórmula Oficial del Anexo VII}}
        \begin{equation*}
          Q_f = \frac{\sum (M_i \cdot PC_i)}{S \cdot 4400} \quad \left[\frac{\kg}{\m^2}\right]
        \end{equation*}
      \end{block}
      \scriptsize
      $\bullet$ $M_i$: Masa inventariada de cada sustancia combustible ($\kg$).\\
      $\bullet$ $PC_i$: Poder Calorífico Inferior de cada material ($\mathrm{kcal/kg}$).\\
      $\bullet$ $S$: Superficie útil de piso del sector de incendio ($\m^2$).\\
      $\bullet$ $4400\,\mathrm{kcal/kg}$: Poder calorífico de la madera patrón de referencia ($\approx 18.4\,\mathrm{MJ/kg}$).
    \end{column}

    \begin{column}{0.48\textwidth}
      \begin{exampleblock}{\textbf{Poderes Caloríficos Comunes}}
        \tiny
        $\bullet$ Madera de referencia: $4400\,\mathrm{kcal/kg}$.\\
        $\bullet$ Papel y cartón: $4000\,\mathrm{kcal/kg}$.\\
        $\bullet$ Plásticos (PVC, nylon): $8000\text{ a }10000\,\mathrm{kcal/kg}$.\\
        $\bullet$ Hidrocarburos y aceites: $10000\text{ a }11500\,\mathrm{kcal/kg}$.\\[3pt]
        \textbf{Significado físico:} Representa cuántos kilos de madera por metro cuadrado equivalen térmicamente a todos los combustibles del recinto.
      \end{exampleblock}
    \end{column}
  \end{columns}
\end{frame}

% --- SLIDE 24: RESISTENCIA F TABLA ---
\begin{frame}{Carga de Fuego y Resistencia Estructural al Fuego ($F$)}
  \small
  El valor resultante de $Q_f$ determina taxativamente la resistencia al fuego mínima de los muros divisorios:

  \vspace{0.2cm}
  \begin{table}[h]
    \centering
    \scriptsize
    \begin{tabular}{p{3.2cm} p{2.2cm} p{2.5cm} p{2.5cm}}
      \toprule
      \textbf{Carga de Fuego ($Q_f$)} & \textbf{Riesgo 2 (Inflam.)} & \textbf{Riesgo 3 (Muy Comb.)} & \textbf{Riesgo 4 (Combustible)} \\
      \midrule
      $Q_f \le 15\,\kg/\m^2$ & $F60$ & $F30$ & --- \\
      $15 < Q_f \le 30\,\kg/\m^2$ & $F90$ & $F60$ & $F30$ \\
      $30 < Q_f \le 60\,\kg/\m^2$ & $F120$ & $F90$ & $F60$ \\
      $60 < Q_f \le 100\,\kg/\m^2$ & $F180$ & $F120$ & $F90$ \\
      $Q_f > 100\,\kg/\m^2$ & $F180$ & $F180$ & $F120$ \\
      \bottomrule
    \end{tabular}
  \end{table}

  \vspace{0.1cm}
  \begin{alertblock}{\textbf{Criterio de Resistencia $F$}}
    \scriptsize El valor numérico ($F30, F60, F120$) indica los minutos durante los cuales el muro o cerramiento conserva su estanqueidad a gases, estabilidad mecánica y aislamiento térmico bajo ensayo de curva de fuego normalizada ISO 834.
  \end{alertblock}
\end{frame}

% --- SLIDE 25: MEDIOS DE ESCAPE Y UAS ---
\begin{frame}{Dimensionamiento de Medios de Escape: Unidades de Ancho (UAS)}
  \begin{columns}[c]
    \begin{column}{0.48\textwidth}
      \begin{block}{\textbf{Definición de UAS}}
        \small
        Espacio libre requerido para que circule una hilera de personas en marcha ordenada:
        \begin{itemize}
          \scriptsize
          \item Primeras $2\,\mathrm{UAS} = 1.10\,\m$ (ancho mínimo de cualquier salida).
          \item Cada unidad adicional suma $+0.45\,\m$.
          \item Ninguna puerta individual medirá $< 0.90\,\m$.
        \end{itemize}
      \end{block}
      \vspace{0.1cm}
      \begin{block}{\textbf{Fórmulas Deterministas}}
        \scriptsize
        Población teórica: $N = \frac{S}{f_o}$ ($f_o$ según uso edilicio)\\[2pt]
        Unidades requeridas: $n = \frac{N}{100}$ (redondeo $\ge 0.5$ arriba)
      \end{block}
    \end{column}

    \begin{column}{0.48\textwidth}
      \centering
      \resizebox{0.9\columnwidth}{!}{%
        \begin{tikzpicture}
          \draw[line width=1.5pt, darkcharcoal] (0,0) rectangle (4,1.8);
          \draw[line width=1pt, dashed, grayborder!80!black] (2,0) -- (2,1.8);
          \draw[<->, >=stealth, line width=1.2pt, firered] (0,2.1) -- (4,2.1) node[midway, above, font=\scriptsize\bfseries] {$2\,\mathrm{UAS} = 1.10\,\m$};
          \node[font=\scriptsize\bfseries] at (1,0.9) {Hilera 1};
          \node[font=\scriptsize\bfseries] at (3,0.9) {Hilera 2};
        \end{tikzpicture}
      }
    \end{column}
  \end{columns}
\end{frame}

% --- SLIDE 26: ERGONOMÍA DE ESCALERAS ---
\begin{frame}{Geometría Ergonómica de Vías y Cajas de Escaleras}
  \begin{columns}[c]
    \begin{column}{0.48\textwidth}
      \begin{block}{\textbf{Fórmula de Paso de Blondel}}
        \begin{equation*}
          2a + p = 0.60\,\m \text{ a } 0.63\,\m
        \end{equation*}
        \scriptsize
        $\bullet$ Alzada ($a$): Altura vertical del escalón ($\le 0.18\,\m$).\\
        $\bullet$ Pedada ($p$): Profundidad horizontal ($\ge 0.26\,\m$).
      \end{block}
      \begin{block}{\textbf{Límites Espaciales de Evacuación}}
        \scriptsize
        $\bullet$ Distancia máxima a salida exterior: $\le 40\,\m$.\\
        $\bullet$ Medios de escape independientes si $n \ge 4\,\mathrm{UAS}$:
        \begin{equation*}
          ME = \frac{n}{4} + 1
        \end{equation*}
      \end{block}
    \end{column}

    \begin{column}{0.48\textwidth}
      \centering
      \resizebox{0.85\columnwidth}{!}{%
        \begin{tikzpicture}[scale=0.85]
          \draw[line width=1.5pt, darkcharcoal] (0,0) -- (1.4,0) -- (1.4,1.0) -- (2.8,1.0) -- (2.8,2.0) -- (4.2,2.0);
          \draw[<->, >=stealth, line width=1pt, firered] (1.4,0) -- (1.4,1.0) node[midway, left, font=\tiny\bfseries] {$a \le 0.18\m$};
          \draw[<->, >=stealth, line width=1pt, utnblue] (1.4,1.0) -- (2.8,1.0) node[midway, below, font=\tiny\bfseries] {$p \ge 0.26\m$};
        \end{tikzpicture}
      }
    \end{column}
  \end{columns}
\end{frame}

% --- SLIDE 27: AEA 90364 SECCIÓN 42 ---
\begin{frame}{Norma AEA 90364 Sección 42: Efectos Térmicos y Cables LS0H}
  \begin{columns}[T]
    \begin{column}{0.48\textwidth}
      \begin{block}{\textbf{Condiciones Ambientales BD2 / BD3 / BD4}}
        \small
        $\bullet$ Locales con alta concentración de personas o riesgo de incendio.\\
        $\bullet$ Prohibición estricta de cables con aislaciones propagadoras de llama o emisoras de humo opaco.\\
        $\bullet$ \textbf{Cables LS0H obligatorios:} Formulados bajo norma \textbf{IRAM 62267} e \textbf{IEC 60332-3}.
      \end{block}
    \end{column}

    \begin{column}{0.48\textwidth}
      \begin{block}{\textbf{Ventajas Químicas del Compuesto LS0H}}
        \small
        $\bullet$ Poliolefinas con hidróxidos metálicos (alúmina trihidratada o hidróxido de magnesio).\\
        $\bullet$ Reacción endotérmica que absorbe calor.\\
        $\bullet$ \textbf{Emisión única de vapor de agua no tóxico.}\\
        $\bullet$ Mantiene la visibilidad en pasillos y anula la corrosión ácida sobre circuitos impresos.
      \end{block}
    \end{column}
  \end{columns}
\end{frame}

% --- SLIDE 28: AEA 90364 SECCIÓN 43 ---
\begin{frame}{Norma AEA 90364 Sección 43: Verificación Térmica Adiabática}
  \begin{columns}[c]
    \begin{column}{0.50\textwidth}
      \begin{block}{\textbf{Inecuación de Energía Pasante}}
        \begin{equation*}
          I_{cc}^2 \cdot t \le k^2 \cdot S^2
        \end{equation*}
      \end{block}
      \scriptsize
      $\bullet$ $I_{cc}$: Corriente eficaz de cortocircuito en el punto considerado ($\mathrm{A}$).\\
      $\bullet$ $t$: Tiempo de despeje de la protección termomagnética o fusible ($\mathrm{s}$).\\
      $\bullet$ $k$: Constante del material ($115\,\mathrm{A\cdot s^{1/2}/mm^2}$ para Cu/PVC; $143$ para Cu/XLPE).\\
      $\bullet$ $S$: Sección transversal nominal del conductor ($\mm^2$).
    \end{column}

    \begin{column}{0.48\textwidth}
      \begin{alertblock}{\textbf{Significado Físico del Límite}}
        \small
        Garantiza que la energía térmica liberada durante el cortocircuito no alcance la temperatura de pirolización del polímero aislante, evitando la ignición del cable antes de la desconexión del interruptor.
      \end{alertblock}
    \end{column}
  \end{columns}
\end{frame}
